% ===================================================================
%  TABLEAU N° 11 — La ville et ses monuments. --- L'incendie.
%
%  Traduction, faite en 2026, en francais standard du Quebec, sur
%  l'ido de texto/io. Regles de la colonne : voir 10-tableau-01.tex.
%  Deux parties, deux numerotations : les cles portent c1 et c2.
%
%  LES HUIT CHOIX PROPRES A CE TABLEAU :
%    (48) wattman ................. conducteur de tramway
%    (64) agents de police ........ policiers
%    (44) employés des postes ..... postiers
%    (33) garçons de banque ....... encaisseurs
%    (50) voiture des ambulances .. ambulance
%    (54) brancard ................ civière
%    (4)  fabriques ............... usines
%         instituteurs ............ enseignants
%
%  CE TABLEAU EST CELUI DES METIERS DE LA VILLE, et c'est pour cela
%  qu'il en compte huit la ou le tableau 10, celui de la marine, n'en
%  comptait que trois. Les noms de fonction sont ce qui vieillit le
%  plus vite dans une langue : le wattman a disparu avec le mot, les
%  garcons de banque avec la tournee des creances a domicile, les
%  employes des postes avec l'administration qui les nommait ainsi.
%  Les choses, elles, restent : le tramway roule toujours, la banque
%  encaisse toujours, la poste distribue toujours.
%
%  CE QUI RESTE, ET POURQUOI. « Hôtel de ville », « palais de
%  justice », « bureau de poste », « caisse d'épargne », « trottoir »
%  sont les mots du Quebec d'aujourd'hui, et les memes qu'en 1926. En
%  revanche « école communale », « préfecture », « préfet » et
%  « gendarmes » designent des institutions FRANCAISES que le Quebec
%  n'a jamais eues : les traduire en institutions d'ici aurait
%  deplace la ville entiere. Ce sont les choses, non la langue.
% ===================================================================

%%K t11-apar-1 apar
\VUsaut{2.0mm}
\VUcentre{13.2pt}{90}{TROISIÈME SÉRIE}
\VUsaut{3.0mm}
\VUfilet{16mm}
\VUsaut{5.6mm}
\VUtitre{15.0pt}{60}{TABLEAU N\textsuperscript{o} 11}
\VUsaut{1.4mm}
\VUpk{11.4pt}{40}{La ville et ses monuments. --- L'incendie.}
\VUsaut{2.2mm}

%%K t11-tit-1 sub
\VUpk{10.2pt}{40}{Le faubourg.}
\VUsaut{3.0mm}

%%K t11-c1-01-1 p
1. --- J'habite une grande ville située à 100 kilomètres de l'Océan et qui est
pourtant un \VUgras{port} \textsuperscript{(1)} de premier ordre. Le fleuve qui
la borde a 400 mètres de largeur et forme une très belle rade.

%%K t11-c1-02-1 p
2. --- Toute la partie de la ville qui se trouve sur les \VUgras{quais}
\textsuperscript{(2)} a un aspect tout maritime et industriel, et forme un
\VUgras{faubourg} \textsuperscript{(3)} qui n'est habité que par des marins et
des ouvriers. Ceux-ci travaillent dans les \VUgras{usines}
\textsuperscript{(4)} et à l'\VUgras{usine à gaz} \textsuperscript{(5)}, dont la
haute \VUgras{cheminée} \textsuperscript{(6)} et le \VUgras{gazomètre} se voient
de très loin.

%%K t11-c1-03-1 p
3. --- Au Moyen Âge, la ville était fortifiée, et l'on trouve encore quelques
restes de ses remparts, en particulier la \VUgras{porte} \textsuperscript{(8)} du
vieux \VUgras{château fort} \textsuperscript{(7)}, flanquée de ses deux
\VUgras{tours} \textsuperscript{(9)}, qui servent aujourd'hui de \VUgras{prison}.

%%K t11-c1-04-1 p
4. --- Dans tout ce \VUgras{quartier}, qui a l'air très vieux, un seul édifice
est relativement moderne : c'est la \VUgras{douane} \textsuperscript{(10)}. Elle
est située entre le quai et la petite \VUgras{rue} \textsuperscript{(11)} qui
conduit au vieil \VUgras{hôtel de ville} \textsuperscript{(12)}. On reconnaît
facilement ce dernier monument grâce au gigantesque \VUgras{beffroi}
\textsuperscript{(13)} qui domine la vieille ville.

%%K t11-c1-05-1 p
5. --- De l'autre côté de la rue \VUgras{s'élève} un édifice construit dans le
même style que la douane : c'est la \VUgras{Bourse} \textsuperscript{(14)}. Une
de ses façades donne sur une petite place où se trouve la \VUgras{préfecture}
\textsuperscript{(15)}. Notre ville, en effet, sans être une capitale, n'en est
pas moins un important chef-lieu de département.

%%K t11-tit-2 sub
\VUsaut{5.0mm}
\VUpk{10.2pt}{40}{Le haut de la ville.}
\VUsaut{3.4mm}

%%K t11-c1-06-1 p
6. --- La garnison, assez nombreuse, est logée dans de vastes \VUgras{casernes}
\textsuperscript{(16)} très salubres; elles s'étendent jusqu'aux grilles du
\VUgras{parc municipal} \textsuperscript{(17)}. Le \VUgras{boulevard}
\textsuperscript{(19)} qui mène à ce parc passe derrière la \VUgras{caisse
d'épargne} \textsuperscript{(18)} et aboutit à la place de la
\VUgras{cathédrale} \textsuperscript{(20)}.

%%K t11-c1-07-1 p
7. --- Tous ces monuments s'étagent en amphithéâtre sur une colline où se trouve
également notre vaste \VUgras{lycée} \textsuperscript{(21)}.

%%K t11-c1-07-2 p
Si l'on descend par une rue un peu en pente, qui rejoint la Grand-Rue, on arrive
au \VUgras{palais de justice} \textsuperscript{(22)}, devant lequel s'étend un
agréable \VUgras{jardin public} \textsuperscript{(26)}. Ce \VUgras{square} n'est
pas très grand, mais il fait au milieu de la ville une oasis de verdure et de
fraîcheur. Aussi est-il très fréquenté par les citadins, qui aiment venir
s'asseoir sous la tente du \VUgras{café} \textsuperscript{(25)} pour écouter les
musiciens militaires qui jouent le jeudi et le dimanche dans le
\VUgras{kiosque} \textsuperscript{(27)} placé entre les pelouses et les massifs.

%%K t11-c1-08-1 p
8. --- Sur une de ces pelouses s'élève une \VUgras{statue} \textsuperscript{(28)}
de bronze, monument érigé à la mémoire d'un de nos concitoyens qui a rendu de
grands services à sa ville natale.

%%K t11-tit-3 sub
\VUsaut{4.0mm}
\VUpk{10.2pt}{40}{Les transports.}
\VUsaut{3.4mm}

%%K t11-c1-09-1 p
9. --- Notre ville n'est pas encore éclairée à l'électricité, elle n'a que des
\VUgras{becs de gaz} \textsuperscript{(29)}. Mais \VUgras{les journaux d'ici} font
campagne pour qu'on améliore ce système d'éclairage, \VUgras{comme} on a déjà
perfectionné \VUgras{le mode de traction des} tramways. \VUgras{Les} vieilles
voitures \VUgras{tirées} par des \VUgras{chevaux} ont été pratiquement
remplacées par de \VUgras{pratiques tramways électriques}
\textsuperscript{(31)}, qui transportent rapidement les voyageurs d'un bout de
\VUgras{la ville à l'autre} pour un prix très \VUgras{bas}.

%%K t11-c1-10-1 p
10. --- Près du palais de justice se trouve la \VUgras{banque}
\textsuperscript{(32)}, où l'on escompte les valeurs commerciales. Les
\VUgras{encaisseurs} \textsuperscript{(33)} vont recouvrer les créances à
domicile.

%%K t11-tit-4 sub
\VUsaut{4.0mm}
\VUpk{10.2pt}{40}{Art et instruction.}
\VUsaut{3.4mm}

%%K t11-c1-11-1 p
11. --- Le bel édifice qui s'élève tout près est le \VUgras{musée}. Il contient
à la fois la \VUgras{bibliothèque} \textsuperscript{(34)} de la ville, les
belles \VUgras{collections d'histoire naturelle} \textsuperscript{(35)} et une
jolie \VUgras{galerie de tableaux} \textsuperscript{(36)} qui constitue une
splendide \VUgras{exposition} permanente.

%%K t11-c1-11-2 p
Il est ouvert au public deux fois par semaine, mais les visiteurs de l'extérieur
peuvent le visiter n'importe quel jour en donnant un pourboire au
\VUgras{gardien} \textsuperscript{(37)}. Les \VUgras{artistes peintres}
\textsuperscript{(38)} viennent y travailler. On les voit \VUgras{devant} leur
chevalet, la \VUgras{palette} à la main, copier les tableaux célèbres.

%%K t11-c1-12-1 p
12. --- Sur la hauteur, derrière le musée, on aperçoit le \VUgras{dôme}
\textsuperscript{(40)} de l'\VUgras{hôpital} \textsuperscript{(39)} et les
bâtiments de l'\VUgras{école normale} \textsuperscript{(41)} où ont été formés
les enseignants de nos écoles communales. Sur la place du Théâtre se trouve un
\VUgras{bureau de poste} \textsuperscript{(43)} installé dans une \VUgras{maison
particulière} \textsuperscript{(42)}.

%%K t11-tit-5 sub
\VUsaut{5.0mm}
\VUpk{11.4pt}{40}{L'incendie.}
\VUsaut{3.0mm}

%%K t11-tit-6 sub
\VUpk{10.2pt}{40}{Au feu!}
\VUsaut{3.4mm}

%%K t11-c2-01-1 p
1. --- Une rumeur terrifiante se répand dans la ville : le feu vient d'éclater
au théâtre! Au moment où j'arrive sur la \VUgras{place} \textsuperscript{(96)},
la foule est déjà massée sur le \VUgras{trottoir} \textsuperscript{(46)} et sur
la \VUgras{chaussée} \textsuperscript{(47)}. Les \VUgras{postiers}
\textsuperscript{(44)} et les \VUgras{facteurs} \textsuperscript{(45)} sortent du
bureau de poste. Un \VUgras{conducteur de tramway} \textsuperscript{(48)} a dû
arrêter sa voiture pour laisser passer l'\VUgras{ambulance}
\textsuperscript{(50)} municipale, qui arrive avec un \VUgras{chirurgien}
\textsuperscript{(52)} et un \VUgras{infirmier} \textsuperscript{(51)} pour
soigner un \VUgras{blessé} \textsuperscript{(53)} qu'on a apporté sur une
\VUgras{civière} \textsuperscript{(54)} auprès du \VUgras{kiosque à journaux}
\textsuperscript{(55)}.

%%K t11-c2-02-1 p
2. --- Une compagnie d'infanterie, qui revenait de l'exercice, s'est arrêtée
pour assurer le service d'ordre avec les \VUgras{gardes municipaux à cheval}
\textsuperscript{(61)}. Mieux que les \VUgras{soldats}
\textsuperscript{(57)} ou que les \VUgras{gendarmes} \textsuperscript{(63)},
\VUgras{les cavaliers}, dont les chevaux se cabrent, font reculer les
\VUgras{curieux} \textsuperscript{(56)}. Les gendarmes, d'ailleurs, sont très occupés. L'un
d'eux indique aux \VUgras{policiers} \textsuperscript{(64)} où l'on doit
transporter une autre \VUgras{victime} \textsuperscript{(65)}, un courageux
pompier tombé d'une échelle.

%%K t11-c2-03-1 p
3. --- L'\VUgras{officier} \textsuperscript{(66)} précise l'endroit où doit se
placer la \VUgras{pompe à vapeur} \textsuperscript{(67)} qui arrive. En
attendant cette puissante machine, il a fait mettre en batterie une
\VUgras{pompe à bras} \textsuperscript{(68)}, dont le long \VUgras{tuyau}
\textsuperscript{(69)} lance l'eau à torrents sur le feu. Six pompiers la
manœuvrent, pendant que leur camarade, qui tient la \VUgras{lance}
\textsuperscript{(70)}, dirige le \VUgras{jet} \textsuperscript{(71)} vers le
centre de l'incendie.

%%K t11-c2-04-1 p
4. --- De l'autre côté du théâtre, on a dressé la grande \VUgras{échelle}
\textsuperscript{(72)}. Elle permet aux pompiers d'approcher des flammes qui
jaillissent au milieu de tourbillons de \VUgras{fumée} \textsuperscript{(75)}
noire.

%%K t11-tit-7 sub
\VUsaut{5.0mm}
\VUpk{10.2pt}{40}{Des actes de courage.}
\VUsaut{3.4mm}

%%K t11-c2-05-1 p
5. --- L'\VUgras{incendie} \textsuperscript{(74)} a pris naissance dans le foyer
du \VUgras{théâtre} \textsuperscript{(73)}, pendant qu'on répétait
l'\VUgras{opéra} dont la première \VUgras{représentation} devait avoir lieu le
lendemain. Heureusement, il n'y avait que peu de \VUgras{spectateurs}
\textsuperscript{(76)} dans la salle. Les malheureux se sont réfugiés sur le
\VUgras{balcon} \textsuperscript{(77)}, où de courageux \VUgras{pompiers}
\textsuperscript{(86)} viennent les secourir avec une échelle et des cordages.

%%K t11-c2-06-1 p
6. --- Sous le \VUgras{péristyle} \textsuperscript{(78)}, où l'\VUgras{affiche}
\textsuperscript{(79)} annonce la représentation, on voit les
\VUgras{musiciens} \textsuperscript{(81)} de l'\VUgras{orchestre}
\textsuperscript{(80)} s'enfuir en emportant leurs instruments :
\VUgras{violon} \textsuperscript{(81)}, \VUgras{flûte} \textsuperscript{(82)},
\VUgras{violoncelle} \textsuperscript{(83)}, \VUgras{trombone}
\textsuperscript{(84)}, \VUgras{tambour} \textsuperscript{(85)}, etc. On
s'efforce de sauver une partie du \VUgras{mobilier} \textsuperscript{(87)}.

%%K t11-c2-07-1 p
7. --- Les \VUgras{acteurs} \textsuperscript{(89)} et les \VUgras{actrices}
\textsuperscript{(88)}, les \VUgras{chanteurs} et les \VUgras{chanteuses} ont dû
fuir avec leur \VUgras{costume} \textsuperscript{(90)} de scène. Le ténor
explique à un \VUgras{journaliste} \textsuperscript{(92)} comment les choses se
sont passées. \VUgras{Le reporter} est accouru dès le premier instant avec les
autorités : le \VUgras{préfet} \textsuperscript{(91)}, le \VUgras{maire}
\textsuperscript{(93)}, le \VUgras{commissaire de police} \textsuperscript{(94)}
et un \VUgras{magistrat} \textsuperscript{(95)}, qui feront enquête pour établir
les responsabilités.
\VUsaut{6.0mm}
\VUfilet{22mm}
